\documentclass[12pt,a4paper]{article}
\usepackage{fontspec}
\setmainfont{Loma}
\XeTeXlinebreaklocale "th"
\XeTeXlinebreakskip=0pt plus 1pt
\usepackage[a4paper,margin=2.2cm]{geometry}
\usepackage{amsmath,booktabs,array,xcolor,hyperref,fancyhdr,setspace}
\definecolor{kimchired}{RGB}{183,51,47}
\definecolor{kimchigreen}{RGB}{66,108,67}
\hypersetup{colorlinks=true,linkcolor=kimchired,urlcolor=kimchigreen}
\pagestyle{fancy}\fancyhf{}\fancyhead[L]{Kimchi AI Learning Project}\fancyfoot[C]{\thepage}
\setstretch{1.2}
\title{\textbf{\color{kimchired}การทำกิมจิ: จากกระบวนการหมักสู่การประยุกต์ใช้ AI}}
\author{เอกสารประกอบการเรียนรู้สำหรับนักเรียน}
\date{}
\begin{document}
\maketitle
\begin{abstract}
เอกสารนี้อธิบายกระบวนการทำกิมจิผักกาดขาว การหมักโดยแบคทีเรียกรดแลกติก แบบจำลองเชิงคณิตศาสตร์ของค่า pH และแนวคิดการประยุกต์ใช้ปัญญาประดิษฐ์เพื่อวิเคราะห์ข้อมูลการหมัก โดยสมการและข้อมูลตัวอย่างใช้เพื่อการเรียนรู้ ไม่ใช่ผลการทดลองจริง
\end{abstract}
\tableofcontents
\newpage
\section{บทนำ}
กิมจิเป็นอาหารหมักแบบดั้งเดิมของเกาหลี โดยกิมจิผักกาดขาวเป็นชนิดที่เป็นที่รู้จักอย่างแพร่หลาย การทำกิมจิเป็นตัวอย่างที่ดีในการเชื่อมโยงวิทยาศาสตร์ คณิตศาสตร์ เทคโนโลยี และปัญญาประดิษฐ์เข้าด้วยกันในโครงงานสำหรับนักเรียน

\section{วัตถุดิบ}
วัตถุดิบพื้นฐานประกอบด้วยผักกาดขาว เกลือ พริกเกาหลี กระเทียม ขิง น้ำปลา ต้นหอม แครอต และหัวไชเท้า ความสะอาดของวัตถุดิบและภาชนะมีความสำคัญต่อคุณภาพของการหมัก

\section{ขั้นตอนการทำ}
\begin{enumerate}
\item เตรียมผักกาดขาว: ล้างและผ่าเป็นส่วนที่เหมาะสม
\item หมักเกลือ: โรยเกลือให้ทั่ว พักจนผักอ่อนตัว แล้วล้างเกลือส่วนเกิน
\item เตรียมซอสกิมจิ: ผสมพริกเกาหลี กระเทียม ขิง น้ำปลา และผักต่าง ๆ
\item คลุกส่วนผสม: ทาซอสให้ทั่วทุกชั้นของผักกาด
\item บรรจุเพื่อหมัก: ใส่ภาชนะสะอาด ปิดฝา และตรวจสอบรสชาติเป็นระยะ
\end{enumerate}

\section{กระบวนการหมัก}
ระหว่างการหมัก แบคทีเรียกรดแลกติกใช้น้ำตาลบางส่วนในวัตถุดิบและสร้างกรดอินทรีย์ ส่งผลให้รสชาติเปรี้ยวขึ้นและค่า pH มีแนวโน้มลดลง ปัจจัยที่มีผลต่อการหมักจริง ได้แก่ อุณหภูมิ ปริมาณเกลือ วัตถุดิบ และชนิดของจุลินทรีย์

\section{แบบจำลองค่า pH}
แบบจำลองเชิงการศึกษาสามารถเขียนเป็น Exponential Decay ได้ดังนี้
\[
\boxed{p(x)=4.2+1.6e^{-0.32x}}
\]
โดย $x$ คือจำนวนวันที่หมัก และ $p(x)$ คือค่า pH โดยประมาณ เมื่อ $x=0$ จะได้ $p(0)=5.8$ และเมื่อเวลาผ่านไปค่าจะค่อย ๆ เข้าใกล้ 4.2

\begin{center}
\begin{tabular}{ccc}
\toprule
วัน & pH โดยประมาณ & แนวโน้ม\\
\midrule
0 & 5.80 & ค่าเริ่มต้น\\
2 & 5.04 & ลดลงเร็ว\\
4 & 4.64 & ความเป็นกรดเพิ่มขึ้น\\
6 & 4.43 & เริ่มลดลงช้าลง\\
8 & 4.32 & เข้าใกล้ค่าคงที่\\
10 & 4.27 & เปลี่ยนแปลงเล็กน้อย\\
12 & 4.23 & ใกล้ 4.2\\
14 & 4.22 & ใกล้ค่าปลายทาง\\
\bottomrule
\end{tabular}
\end{center}

\section{การประยุกต์ใช้ AI}
เมื่อมีข้อมูลจากการทดลองจริง สามารถนำ AI มาช่วยทำนายค่า pH จากตัวแปร เช่น วันหมัก อุณหภูมิ ปริมาณเกลือ และความชื้น รวมถึงใช้ Computer Vision วิเคราะห์สีและลักษณะของกิมจิ เพื่อสนับสนุนการตัดสินใจว่าจะหมักต่อหรือควรนำไปแช่เย็น

\section{จาก Zero-shot Prompt สู่ Few-shot Prompt}
Zero-shot Prompt คือการสั่งงาน AI โดยให้โจทย์โดยตรงโดยไม่มีตัวอย่างผลลัพธ์ ขณะที่ Few-shot Prompt เพิ่มตัวอย่างที่ต้องการ 2--3 ตัวอย่าง ทำให้ AI เข้าใจรูปแบบ ภาษา โครงสร้าง และระดับรายละเอียดได้ชัดขึ้น สำหรับโครงงานนี้ Few-shot Prompt ช่วยให้การสร้างภาพ อินโฟกราฟิก คำอธิบาย และเนื้อหาสำหรับนักเรียนมีความสม่ำเสมอมากกว่าเดิม

\section{บทสรุป}
โครงงานกิมจิสามารถใช้เป็นกิจกรรมบูรณาการได้ตั้งแต่การทำอาหารและศึกษากระบวนการหมัก ไปจนถึงการสร้างกราฟ แบบจำลองทางคณิตศาสตร์ การออกแบบสื่อ และแนวคิด AI สิ่งสำคัญคือแยกให้ชัดเจนระหว่างข้อมูลทดลองจริงกับข้อมูลจำลองเพื่อการเรียนรู้
\end{document}
